\section{AI Use Log}
\label{app:ailog}

% ---------------------------------------------------------------------
% Layout: each task starts with a small table (tool + prompt), followed
% by short paragraphs. Each task is separate: edit only your own task.
% \aiprompt{Tool, version and date}{Prompt, pasted exactly}
% ---------------------------------------------------------------------
\newcommand{\aiprompt}[2]{%
  \par\noindent
  \begin{tabularx}{\linewidth}{@{}>{\bfseries}p{2.2cm}Y@{}}
    \toprule
    Tool   & #1 \\
    Prompt & \itshape #2 \\
    \bottomrule
  \end{tabularx}\par\medskip}

% =====================================================================
\subsection{Task 1: Business Ecosystems}

\aiprompt{GPT-5.6 Luna}{I am analyzing S3NS, a partnership between Thales and
Google Cloud that provides trusted cloud for French public-sector and
critical-infrastructure organistations. Identify the stakeholders, state each
one's role, and describe the key relationships and dependencies between them.
Keep the answer under 400 words.}

\paragraph{Claims checked.}
The core actors and roles are correct \parencite{s3ns2025premi3ns}.
``Majority shareholder'' is imprecise: S3NS is fully owned by Thales, with no
disclosed Google equity.

\paragraph{Errors found and corrected.}
\textit{(1) Single-firm view.} The AI centred the ecosystem on S3NS as a hub.
Using ecosystem-as-structure, we start from the alliance's value proposition
instead \parencite{adner2017}\parencite[slides~8--13]{windekilde2026ecosystem}.

\textit{(2) Dependency shown as balanced.} S3NS needs Google's technology, but
Google does not need S3NS. This is a unilateral complementarity \parencite[p.~114]{overby2021}, and the
technology is co-specialised \parencite[slide~30]{windekilde2026ecosystem}. Every Google update must also be validated by S3NS
before deployment \parencite{s3ns2025premi3ns}.

\textit{(3) Wrong motive for Google.} Google gains access to a market that
requires immunity from non-European extraterritorial law, not ``security
capabilities'' \parencite{s3ns2025premi3ns}. The French state was added as a
stakeholder.

\textit{(4) No indirect actors.} It listed no indirect actors such as the extraterritorial law.


% =====================================================================
\subsection{Task 2: Network Effects and Path Dependency}

\aiprompt{\fillin{Tool, version, date}}{\fillin{Paste the prompt exactly}}

\paragraph{Claims checked.}
\fillin{What was checked, and against which chapter, lecture or case source}

\paragraph{Errors found and corrected.}
\fillin{For each error: the flaw, the course concept that exposes it, the
evidence, and how the main text changed}

% =====================================================================
\subsection{Task 3: Multisided Platform Mapping}

\aiprompt{Grok 4.5, 07/10/2026}{S3NS is a trusted cloud alliance between Thales and Google Cloud, operating under the brand PREMI3NS for French public-sector and regulated private-sector clients. Treat PREMI3NS as a multisided platform. List the distinct sides that interact through it, explain the cross-side and same-side network effects between them, and propose a pricing model: which side should be charged and which side should be subsidised, and why. Limit your answer to 400 words.}

\paragraph{Claims checked.}
The answer invents no prices and leaves the premium unquantified, as the price list does \parencite{s3nsprices}. We checked it against the price list, the course and the case sources. Three claims hold: usage-based pricing of clients matches the per-unit price list, the client and ISV sides exist \parencite{s3nsecosystem}, and the effect from clients to ISVs and the weak same-side effect among clients match Table~\ref{tab:effects}.

\paragraph{Errors found and corrected.}
\textit{(1) Platform owners treated as sides.} The AI lists Google (supply side) and Thales/S3NS (operator side) next to clients and ISVs. Sides are user groups that interact through the platform \parencite[p.~150]{overby2021}. Google is a supplier, as in the reseller model \parencite[slide 40]{windekilde2026msp}, and S3NS is the platform itself. We kept two sides, clients and complementors (Figure~\ref{fig:msp}).

\textit{(2) Subsidy logic contradicts MSP theory.} It subsidises Google and Thales, the alliance partners, although a subsidy belongs to a price-sensitive user side that creates value for the other side \parencite{rochet2006}, as its own final sentence says. It leaves ISV pricing open (``charged \ldots{} or subsidised'') although the ISV side is small and multi-homing \parencite{futurum2026}, and it does not say who funds the subsidy. We subsidise the ISVs, show that Thales funds this as anchor client in the SAP case \parencite{thalessap2026}, and expect fees later \parencite{armstrong2006}.

\textit{(3) Concepts misclassified.} More users raising the value of ``Google's innovation pipeline'' is supply-side economies of scale, not a network effect \parencite[p.~124]{overby2021}, \parencite[slide 5]{windekilde2026msp}, and the scale belongs to Google (Section~\ref{sec:ValueGrowth}). The ``complementary effects'' among Google, Thales and ISVs are complementarities, unilateral between Google and S3NS \parencite[p.~114]{overby2021}, not same-side effects. Among ISVs the same-side effect is competition (Table~\ref{tab:effects}).

% =====================================================================
\subsection{Task 4: Cybersecurity Economics}

\aiprompt{\fillin{ChatGPT, 6.1, 04/10/2026}}{\fillin{S3NS is a trusted cloud alliance between Thales and Google Cloud. S3NS is a French company fully controlled by Thales, and it operates the PREMI3NS cloud service for French public-sector and regulated private-sector clients, using licensed Google Cloud technology. PREMI3NS holds the SecNumCloud qualification from ANSSI. Software vendors and integration partners also run services on the platform. Treat platform security as a good provided jointly by the alliance partners. Explain how cybersecurity responsibility and investment should be allocated between Thales/S3NS and Google Cloud, and which parts of the attack surface each partner controls. Identify who would bear the losses if the platform were breached. Finally, assess whether sovereign domestic hosting in France makes the platform more secure. Limit your answer to 400 words.}}

\paragraph{Claims checked.}
There is no invented numbers in this AI output. Below is claimed which have been checked and how we evaluate them.

That S3NS exclusively operates in France and that there is quarantines, ananlyses and validations of the technology before it gets deployed. A claim that is correctly stated according to \parencite{s3ns2025premi3ns}. Additionally, the AI claims that ANSSI warns about that the qualification of a cloud offer dont establish the security of applications hosting on it. A claim that is confirmed by ANSSI FAQ which says that the qualification itself does not guarantee clients service level \parencite{anssi_faq}. 

\paragraph{Errors found and corrected.}
\fillin{For each error: the flaw, the course concept that exposes it, the
evidence, and how the main text changed}

% =====================================================================
\subsection{Task 5: Business Model Structure}

\aiprompt{\fillin{Tool, version, date}}{\fillin{Paste the prompt exactly}}

\paragraph{Claims checked.}
\fillin{What was checked, and against which chapter, lecture or case source}

\paragraph{Errors found and corrected.}
\fillin{For each error: the flaw, the course concept that exposes it, the
evidence, and how the main text changed}

